%===============================================================
% Apéndice 1: Implicaciones ambientales y de desarrollo sostenible
% Proyecto: ModuTEC – Modulación y Demodulación de Señales en
%           Sistemas Embebidos
% TFG CE-5601 | Instituto Tecnológico de Costa Rica
%===============================================================

\appendix
\chapter{Implicaciones ambientales y de desarrollo sostenible}
\label{apx:sostenibilidad}

El presente apéndice analiza las implicaciones del proyecto ModuTEC
desde una perspectiva amplia de sostenibilidad, considerando su impacto
actual y potencial sobre la sociedad, el entorno y el ejercicio
profesional de la ingeniería en computadores. El proyecto consistió en
el desarrollo de un sistema embebido para la modulación y demodulación
de señales AM y ASK/OOK sobre hardware de bajo costo, integrando
procesamiento digital de señales, visualización en tiempo real y
evaluación cuantitativa del desempeño. Si bien su alcance es académico,
las decisiones de ingeniería adoptadas a lo largo del desarrollo tienen
implicaciones que trascienden el entorno controlado en el que fue
concebido.

%---------------------------------------------------------------
\section{Identificación y análisis de los impactos asociados al proyecto}
%---------------------------------------------------------------

El impacto social más directo del proyecto se deriva de la accesibilidad
económica del hardware seleccionado. La dependencia de simuladores
especializados o equipos de laboratorio de elevado costo limita las
posibilidades de experimentación práctica en entornos educativos con
recursos restringidos, profundizando brechas entre instituciones con
distintos niveles de infraestructura. Al implementar la solución sobre
un microcontrolador de bajo costo, el proyecto demuestra que es posible
abordar técnicas de procesamiento digital de señales sin depender de
equipamiento especializado, lo que amplía el acceso a herramientas de
formación práctica en áreas como las comunicaciones digitales y los
sistemas embebidos. Esta característica tiene un impacto social concreto
sobre la equidad en la educación técnica, tanto en el contexto
institucional del Instituto Tecnológico de Costa Rica como en otros
entornos formativos con limitaciones similares.

Desde el punto de vista económico, la decisión de diseño que prioriza
el bajo costo no es únicamente una restricción presupuestaria, sino una
postura de ingeniería con implicaciones sobre la sostenibilidad del
sistema. Al demostrar que la complejidad técnica no exige inversiones
elevadas en infraestructura, el proyecto aporta evidencia práctica
que puede incidir sobre las decisiones de equipamiento en laboratorios
de formación, reduciendo costos operativos sin comprometer la calidad
de la experiencia educativa. Esta lógica tiene valor económico acumulado
cuando se replica en múltiples instituciones y contextos.

En cuanto al aspecto ambiental, el hardware utilizado presenta un
consumo energético reducido en comparación con plataformas de mayor
capacidad computacional, lo que representa una ventaja durante la
operación continua del sistema. La arquitectura modular adoptada, que
permite extender la funcionalidad del sistema sin reemplazar los
componentes existentes, reduce la generación de residuos electrónicos
asociados a la obsolescencia tecnológica. No obstante, como todo
hardware electrónico, los componentes del sistema generan residuos al
final de su vida útil cuya gestión inadecuada puede tener efectos
negativos sobre el entorno. Esta consideración es relevante para el
ejercicio responsable de la ingeniería y recae tanto en los usuarios
finales como en los sistemas institucionales de gestión de residuos
tecnológicos.

Respecto al aspecto legal, el proyecto fue desarrollado íntegramente
mediante herramientas de software de código abierto cuyos términos
de licencia permiten su uso académico y de investigación sin
restricciones significativas. El sistema opera mediante comunicación
serial USB en un entorno completamente cerrado, sin emisión de señales
hacia el medio externo, por lo que no genera afectaciones en materia
de regulación del espectro radioeléctrico bajo su configuración actual.
Sin embargo, extensiones futuras del sistema hacia comunicaciones
inalámbricas requerirían una evaluación regulatoria cuidadosa, aspecto
que el ingeniero responsable del desarrollo deberá contemplar
explícitamente en etapas posteriores.

Finalmente, en términos de seguridad y salud, el sistema opera con
niveles de tensión reducidos propios de la alimentación USB, lo que
minimiza el riesgo de accidentes eléctricos durante su manipulación
en entornos de laboratorio. Desde la perspectiva de la seguridad
informática, la ausencia de mecanismos de cifrado en la comunicación
serial no representa un riesgo bajo el alcance académico del proyecto,
pero constituye una limitación a considerar en aplicaciones futuras
con requerimientos de integridad o confidencialidad de datos. En
materia de salud, el uso prolongado de la interfaz gráfica de monitoreo
introduce factores ergonómicos comunes al trabajo con dispositivos de
visualización, cuya gestión recae en las prácticas de higiene postural
y de descanso recomendadas para estos entornos.

%---------------------------------------------------------------
\section{Relación del proyecto con los Objetivos de Desarrollo Sostenible}
%---------------------------------------------------------------

El proyecto ModuTEC guarda relación directa con varios de los Objetivos
de Desarrollo Sostenible establecidos por las Naciones Unidas en la
Agenda 2030. El vínculo más evidente se establece con el ODS~4, orientado
a garantizar una educación de calidad e inclusiva, dado que la solución
desarrollada facilita el acceso a herramientas de experimentación práctica
en telecomunicaciones y sistemas embebidos sin depender de infraestructura
costosa. Esta contribución se concreta en la posibilidad de que entornos
educativos con recursos limitados incorporen prácticas formativas de mayor
profundidad técnica a partir de plataformas replicables y accesibles.

En el plano de la innovación tecnológica, el proyecto se alinea con el
ODS~9, cuya meta de fomentar la investigación científica y el desarrollo
de infraestructura tecnológica accesible encuentra en esta solución un
ejemplo concreto de cómo la ingeniería aplicada puede generar conocimiento
técnico reproducible sobre plataformas de bajo costo. La arquitectura
modular y extensible del sistema refuerza este alineamiento al proporcionar
una base reutilizable para desarrollos futuros.

La accesibilidad económica del hardware seleccionado tiene implicaciones
directas sobre el ODS~10, orientado a la reducción de desigualdades. Al
eliminar la barrera de costo como restricción de acceso a herramientas
de formación práctica en ingeniería, el proyecto contribuye a reducir la
brecha tecnológica entre instituciones con distintos niveles de recursos,
favoreciendo una distribución más equitativa de las capacidades de
experimentación técnica. Esta dimensión adquiere especial relevancia en
el contexto latinoamericano, donde la disparidad en infraestructura
educativa es una limitación real para el desarrollo de competencias
técnicas avanzadas.

Desde la perspectiva ambiental, el diseño del sistema con criterios de
bajo consumo energético y modularidad extensible se relaciona al ODS~12, orientado a la producción y el consumo responsables. La
posibilidad de extender la funcionalidad del sistema sin reemplazar el
hardware existente favorece un modelo de uso tecnológico más sostenible,
reduciendo la generación de residuos electrónicos asociados a ciclos de
reemplazo acelerado. Finalmente, el desarrollo del proyecto en el marco
de una institución pública de educación superior orientada al impacto
social refleja el espíritu colaborativo del ODS~17, en cuanto que el
conocimiento generado tiene potencial de ser transferido y adoptado por
otros grupos académicos o instituciones con necesidades similares.

%---------------------------------------------------------------
\section{Proyecciones futuras e impacto potencial del proyecto}
%---------------------------------------------------------------

A mediano y largo plazo, el mayor impacto potencial del proyecto se
deriva de su capacidad de ser replicado, extendido y adoptado como
recurso educativo en múltiples contextos. La disponibilidad de una
plataforma funcional y documentada para la experimentación con técnicas
de modulación y demodulación puede incidir positivamente en la calidad
de la formación práctica de estudiantes de ingeniería en Costa Rica y
en otros países de la región, especialmente en instituciones donde el
acceso a equipamiento especializado es limitado. En términos sociales,
esta proyección apunta hacia una democratización progresiva del
conocimiento aplicado en comunicaciones digitales, con impacto sobre
la equidad educativa en el área tecnológica.

Desde la perspectiva tecnológica y ambiental, el proyecto se inserta
en una tendencia más amplia hacia la computación de borde, en la que
el procesamiento se desplaza desde infraestructuras centralizadas hacia
dispositivos embebidos de bajo consumo. Esta tendencia tiene
implicaciones ambientales favorables en términos de eficiencia energética
agregada, y posiciona soluciones como la desarrollada en este proyecto
como antecedentes relevantes para el diseño de sistemas distribuidos
de bajo impacto energético.

Económicamente, la formación de ingenieros con experiencia práctica en
plataformas de bajo costo fortalece el ecosistema local de desarrollo
tecnológico, ampliando el perfil de soluciones abordables sin depender
de importación de tecnología especializada. La consolidación de líneas
de investigación aplicada en sistemas embebidos dentro del Instituto
Tecnológico de Costa Rica puede derivar en mayor producción académica
y en proyectos de vinculación con el sector productivo, generando
impacto económico que trasciende el alcance individual de este trabajo.

En materia legal, la eventual extensión del sistema hacia comunicaciones
inalámbricas demandará una alineación explícita con el marco regulatorio
aplicable en materia de uso del espectro radioeléctrico, aspecto que
el ingeniero responsable de ese desarrollo deberá contemplar desde
etapas tempranas del diseño. En cuanto a seguridad y salud, la proyección
del área hacia aplicaciones de monitoreo en entornos críticos, como
salud o gestión ambiental, elevaría considerablemente los requisitos
de confiabilidad, seguridad informática y ergonomía exigibles a sistemas
de esta naturaleza. La arquitectura modular y los principios de evaluación
cuantitativa del desempeño incorporados en este proyecto son compatibles
con esos estándares más exigentes, constituyendo una base técnica sobre
la cual pueden desarrollarse soluciones de mayor complejidad en el futuro.
